% ===========================================================================
\section{The dictionary}\label{sec:dict}
\warmth{0}\quad\whisper{Hard in one domain, easy in the other. The headline: convolution becomes multiplication.}

\begin{strand}
The transform earns its fame through a \keyword{dictionary}: an awkward operation on $f$ turns
into a simple one on $\widehat f$. The crown jewel is the \keyword{convolution theorem} —
smoothing, filtering, and blurring are all convolutions, and in frequency they are just
pointwise products.
\end{strand}

\block{The translation table}
\noindent\begin{tabularx}{\linewidth}{@{}l l L@{}}
\toprule
{\color{indigo}on $f$} & {\color{indigo}on $\widehat f$} & \multicolumn{1}{l}{\color{indigo}name / use}\\
\midrule
$f(x-a)$              & $e^{-2\pi i a\xi}\,\widehat f(\xi)$  & shift $\to$ phase\\
$e^{2\pi i a x}f(x)$  & $\widehat f(\xi-a)$                  & modulation $\to$ shift\\
$(f*g)(x)$            & $\widehat f(\xi)\,\widehat g(\xi)$   & \keyword{convolution $\to$ product}\\
$f(x)\,g(x)$          & $(\widehat f*\widehat g)(\xi)$       & product $\to$ convolution\\
$f'(x)$               & $2\pi i\,\xi\,\widehat f(\xi)$       & derivative $\to$ multiply by $\xi$\\
$\overline{f(-x)}$    & $\overline{\widehat f(\xi)}$         & symmetry / realness\\
\bottomrule
\end{tabularx}

\block{The two theorems that power the table}

\begin{theorem}[convolution theorem]\label{thm:conv}
With $(f*g)(x)=\int f(y)\,g(x-y)\dd y$, one has $\widehat{f*g}=\widehat f\cdot\widehat g$.
\end{theorem}
\begin{proof}
Write the transform of $f*g$ as a double integral and substitute $u=x-y$:
\[
  \widehat{f*g}(\xi)=\int\!\!\int f(y)g(x-y)e^{-2\pi i\xi x}\dd y\dd x .
\]
\TODO{split $e^{-2\pi i\xi x}=e^{-2\pi i\xi y}e^{-2\pi i\xi(x-y)}$ and factor the double integral into $\widehat f(\xi)\,\widehat g(\xi)$}
\end{proof}

\begin{theorem}[Parseval / Plancherel]\label{thm:parseval}
The transform preserves inner products and energy:
$\inner{f}{g}=\inner{\widehat f}{\widehat g}$, and in particular $\norm{f}_2=\norm{\widehat f}_2$.
\end{theorem}

\begin{yourturn}
Use the dictionary, not integrals.
\begin{itemize}[leftmargin=1.3em,itemsep=1pt]
  \item To \keyword{filter} a signal (blur it with a kernel $g$), you convolve $f*g$. In frequency
        that is just $\widehat f\cdot\widehat g$, so a filter is a \fillin[3cm] of the spectrum.
  \item Solving $f''=h$: transforming turns it into $(2\pi i\xi)^2\widehat f=\widehat h$, so
        $\widehat f(\xi)=\fillin[2.5cm]$. A differential equation became \emph{division}.
\end{itemize}
\recall{A filter \emph{reshapes} (multiplies) the spectrum; and $\widehat f=\widehat h/(-4\pi^2\xi^2)$. Calculus $\to$ algebra is the whole point.}
\end{yourturn}

\begin{strand}
The discrete mirror: convolution of length-$N$ vectors is diagonalized by the DFT, i.e.\ every
\keyword{circulant matrix} is diagonalized by the Fourier matrix, with the transforms as its
eigenvalues. Shift-invariance $=$ circulant $=$ diagonal-in-Fourier; one fact, three faces.
\end{strand}
